% Guillaumet trip report.
%
% Two targets from one source:
%   - PDF:  pdflatex guillaumet.tex   (everything in the preamble applies)
%   - Blog: ./build.sh                (Pandoc reads \title/\date and the body;
%                                      it ignores the typography packages, so
%                                      the HTML falls back to the site stylesheet)
%
% Keep body markup to what Pandoc understands: \section, \emph, \textbf,
% itemize, tabular, \footnote, \href. Don't \renewcommand anything used in the
% body; Pandoc expands user macros and the expansion leaks into the HTML.

\documentclass[11pt]{article}

% ---------------------------------------------------------------------------
% Page, type, spacing (PDF only)
% ---------------------------------------------------------------------------
\usepackage[a5paper,margin=18mm,top=22mm,bottom=24mm]{geometry}
\usepackage[T1]{fontenc}
\usepackage[utf8]{inputenc}
\usepackage[osf]{libertine}     % Linux Libertine, old-style figures
\usepackage{microtype}
\usepackage{setspace}
\setstretch{1.12}
\usepackage{parskip}
\setlength{\parskip}{0.55\baselineskip}

\usepackage{xcolor}
\definecolor{granite}{HTML}{3B3F45}
\definecolor{glacier}{HTML}{4A7B9D}
\definecolor{rule}{HTML}{B8BEC6}

\usepackage{titlesec}
% Title and rule in one unbreakable box, so a page break cannot split them.
\newcommand{\ruledtitle}[1]{\parbox[t]{\linewidth}{#1\par\vspace{2pt}{\color{rule}\hrule height 0.4pt}}}
\titleformat{\section}
  {\color{granite}\normalfont\Large\itshape}
  {}{0pt}{\ruledtitle}
\titlespacing*{\section}{0pt}{1.6\baselineskip}{0.6\baselineskip}

\usepackage{booktabs}
\usepackage{graphicx}
\usepackage[font={small,it},labelformat=empty,skip=5pt]{caption}
\setlength{\textfloatsep}{1.2\baselineskip}
\setlength{\intextsep}{1.2\baselineskip}

% Photo layouts (wrapphoto, photopair, banner) live in posts/photos.sty.
\usepackage{photos}
\setlength{\bannerheight}{75mm}
\renewcommand{\bannerfocus}{0.3}
\usepackage{array}
\usepackage{lettrine}
\usepackage{fancyhdr}
\pagestyle{fancy}
\fancyhf{}
\renewcommand{\headrulewidth}{0pt}
\fancyfoot[C]{\small\color{granite}\thepage}
\fancyhead[R]{\small\color{granite}\itshape Guillaumet}

\usepackage[hidelinks]{hyperref}
\hypersetup{pdftitle={Guillaumet},pdfsubject={Patagonia trip report}}

% ---------------------------------------------------------------------------
% Title block (PDF only; Pandoc ignores titling hooks and uses \title/\date)
% ---------------------------------------------------------------------------
\usepackage{titling}
\setlength{\droptitle}{-1.5\baselineskip}
\pretitle{\begin{center}
  {\color{glacier}\footnotesize\scshape patagonia\par}
 \vspace{0.6\baselineskip}
  \color{granite}\Huge\itshape}
\posttitle{\par\vspace{0.5\baselineskip}
  {\color{granite}\small Aguja Guillaumet, 2580\,m \textperiodcentered{} Fitz Roy massif\par}
  \end{center}}
\preauthor{}\postauthor{}
\predate{\vspace{-0.6\baselineskip}\begin{center}\color{granite}\small}
\postdate{\par
  \vspace{0.3\baselineskip}
  {\color{rule}\rule{0.35\linewidth}{0.4pt}\par}
  \end{center}\vspace{0.5\baselineskip}}

\title{Guillaumet}
\author{}
\date{January 2nd - 3rd, 2026}

% ---------------------------------------------------------------------------
\begin{document}
\begin{banner}
  \includegraphics[width=\linewidth]{img/guillaumet/banner.jpeg}
\end{banner}
\maketitle

\lettrine[lines=3,findent=2pt,nindent=0pt]{T}{he main goal} of my trip to
Patagonia was to climb Guillaumet, a small peak on the far north side of the
Fitz Roy massif with an elevation of 2580 meters. This would be my first real
alpine climb and the first time I worked with a climbing guide.

\section{Planning}

The planning started a few months earlier in the US. After finding a guiding
company that took clients up Guillaumet, I reached out and explained my
climbing experience. They sounded confident that I had enough experience and
agreed to take me on as a client. I really wanted to do this climb with a
friend, but not having a partner who could take the time off to make it down to
Argentina forced my hand, and I was happy to take on this challenge with
someone very experienced in this type of terrain.

Packing went smoothly. I only needed to purchase a new pair of winter boots,
since my ice climbing boots were a bit too stiff to comfortably walk in for a
substantial amount of time. I had the rest of the gear, including stuff for
winter camping, rock climbing and ice climbing. I started to work on my cardio
a bit more and got out on the snow a few times to try and get a bit more
comfortable on that type of terrain. Ideally I would have done a little more
training with a heavy pack, but this was something I only realized later.
Feeling as ready as I could be, I left the US for Argentina.

\section{El Chalt\'en}

Shortly after arriving in El Chalt\'en, I met with my guide Ramiro. He seemed
like a cool guy and we wasted no time starting to plan our climb. The first day
I arrived was the start of a great weather window. We considered leaving for
the mountains the very next day, but due to my jet lag and an optimistic
forecast for the next week, we decided to hold off for a few days. This would
give us some time to climb together around town and get to know each other
better. It would also give him a bit more confidence that I truly knew what I
was doing.

The following day we walked about 5 minutes outside of town to the local
crag. We climbed up a fun 5-pitch sport route that went through a large roof.
We then rappelled down and walked a few minutes to another route, which we did
in 3 pitches. The skies were perfectly clear and I had my best views of the
Fitz Roy massif so far. Stunning! After the second climb we hiked down the
back side of the crag and parted ways, planning to start our hike into the
mountains in two days.

I spent the next day packing, getting food and generally resting. I was
feeling really good and psyched to start the main objective.

\section{The approach}

\begin{wrapphoto}{l}
  \includegraphics[width=\linewidth]{img/guillaumet/approach_1.jpg}
  \caption{Following the Rio Electrico toward Piedra del Fraile.}
\end{wrapphoto}
The following day, January 2nd, we met around 10:30\,am to drive over to the
Electrico valley and start our hike in. We were still a bit unsure of the
weather forecast. High winds were forecast for the following morning, which
was nerve-racking. We considered changing plans to climb Aguja de l'S, but
ultimately decided to go for the original objective. This turned out to be the
correct decision. After a quick gear check we made the short drive and started
hiking in. It was all finally happening!

The hike starts by following Rio Electrico for a few miles before reaching
Piedra del Fraile. This is a magical little campground with a nice host
family. After stopping here for a short rest we continued on. From there the
hike got much more strenuous. It gains almost 3500\,ft of elevation over about
2 miles. With heavy packs this felt tough, but we moved consistently and
eventually reached our camp for the night, Piedra Negra. This took about 4
hours in total.

\begin{photopair}
  \includegraphics[width=0.49\linewidth]{img/guillaumet/approach_2.jpg}\hfill
  \includegraphics[width=0.49\linewidth]{img/guillaumet/approach_3.jpg}
  \caption{Lupines and the refugio at Piedra del Fraile. The grind up to Piedra Negra, the north end of the massif coming into view.}
\end{photopair}

\begin{wrapphoto}{r}
  \includegraphics[width=\linewidth]{img/guillaumet/approach_4.jpg}
  \caption{Camp at Piedra Negra, tucked behind a wind wall. Guillaumet in the background.}
\end{wrapphoto}
After arriving at camp we set up the tent, got water, and settled in for the
night. We were able to get updated weather forecasts from Piedra del Fraile
via radio, as well as an inReach message from Ramiro's girlfriend. Despite the
wind when we arrived at camp, the forecast looked good and the wind
consistently died down through the night. We had some dehydrated meals and an
empanada for dinner, along with some chicken broth to drink. Temperatures were
quite comfortable, though I am not sure of the exact numbers.

\section{Summit day}

At 3:30\,am the alarm went off. We boiled some water for instant coffee and
put on our boots and jackets. By 4\,am we had started the hike to the base of
the route. The hike involved 45 to 60 minutes of easy talus hiking. After this
we arrived at a snowfield and put crampons on. From here there was about an
hour of moderate snow, with a small section of steep snow. All of this was
done un-roped and felt pretty comfortable. I was surprised at how good I felt
physically after the hike the previous day and limited sleep. I did not feel
sore at all and was really psyched! Pretty minimal nerves, which was also
surprising. After roughly 2 hours of hiking, we arrived at the base of the
route.

\begin{figure}[H]
  \centering
  \includegraphics[width=\linewidth]{img/guillaumet/climb_1.jpg}
  \caption{Alpenglow from the base of the route, looking back down the approach.}
\end{figure}

We took 30 minutes or so to transition to rock climbing gear. Crampons and
boots in our packs, harnesses and climbing shoes on, and ropes flaked. By this
point it was quite bright, so we also stowed our headlamps and enjoyed the
last of the alpenglow as the sun was rising. And then Ramiro was off leading
the first pitch.

The first several pitches were quite moderate. I was delighted to find the
climbing really similar to Little Cottonwood Canyon: pristine white granite
with lots of splitter cracks for gear and jams. I felt really comfortable on
the terrain and climbed quickly. The only slowdown on this portion of the
climb was when Ramiro dislodged a rock in a chimney and core-shot one of the
ropes. I didn't know it at the time from my belay stance, but he had to cut 4
to 5 feet off the end of the rope. When I got to the same spot I opted to
climb the slab on the outside of the chimney to avoid a similar issue. After 5
or so pitches we reached the crux of the route, a 5.10 splitter in a corner!

The crux was probably the best part of the route. While it was strenuous, the
jams were great and I made short work of it without falling. The crack
consisted of many sizes, from 0.3 to 3, fortunately nothing too big. After the
crux we were roughly halfway done with the climbing. An airy traverse pitch
was the next highlight, moving across a juggy flake with smears for feet. It
felt quite easy, but was cool to have some exposure way up in the alpine.

% Mixed portrait + landscape row: widths chosen so heights match
% (0.353 * 4/3 = 0.627 * 3/4).
\begin{photopair}
  \includegraphics[width=0.353\linewidth]{img/guillaumet/climb_2.jpg}\hfill
  \includegraphics[width=0.627\linewidth]{img/guillaumet/climb_3.jpg}
  \caption{Several pitches up the route, Lago Electrico and the Marconi glacier far below. Pristine white granite, with the summit snowfield above.}
\end{photopair}

After 4 to 5 more pitches we arrived at the base of the summit snowfield. We
switched back to boots and crampons before short-roping up this final pitch to
the summit. My ego wasn't thrilled with being short-roped up this fairly
moderate snowfield, but not unexpected on a guided climb. At 10:30\,am we
arrived at the summit!

\begin{photopair}
  \includegraphics[width=0.49\linewidth]{img/guillaumet/summit_1.jpg}\hfill
  \includegraphics[width=0.49\linewidth]{img/guillaumet/summit_2.jpg}
  \caption{Summit selfie with Ramiro. Taking a break among the summit blocks, the ice cap behind.}
\end{photopair}

\section{Descent}

After hanging out on the summit for a while, taking photos, eating snacks, and
chatting with an American couple who was also on the summit, we started the
descent.

The first part consisted of walking down the summit snow slope, back to the end
of the rock section. Again, I was short-roped down this section before arriving
at the first of 5 rappels. The raps consisted of 3 fairly straight forward rock
rappels to get back to the top of Amy's couloir. Then two more rappels down the
couloir itself. The first section was pretty standard and similar to a lof of
descents I have done in the past. The only difference was how much it traversed.
This made for some interesting rope pulls but fortunately nothing got snagged
and we were quickly at the top of Amy's.

\begin{photopair}
  \includegraphics[width=0.49\linewidth]{img/guillaumet/descent_1.jpeg}\hfill
  \includegraphics[width=0.49\linewidth]{img/guillaumet/descent_3.jpeg}
  \caption{Tat and rings at two of the rock rappel anchors.}
\end{photopair}

The second section of rappels was quite intimidating for me. We first went down
the narrow top section of the couloir which was starting to melt out
significantly. This made for a lot of rock fall, and since I was going down
second, I had to be very careful to not kick anything down on Ramiro. The second
rappel in this section was initially much less scary since the couloir had
opened up, though rockfall was still a serious concern. However at the bottom of
this rappel there was a large bergshrund. This looked almost un-passable to me,
but Ramiro calmly explained that we would be rappeling into the shrund and then
climbing out the other side. This ended up being a really cool experience, but
quite scary at the time since I had never been that close to a crevasse.

\begin{photopair}
  \includegraphics[width=0.49\linewidth]{img/guillaumet/descent_2.jpeg}\hfill
  \includegraphics[width=0.49\linewidth]{img/guillaumet/descent_4.jpeg}
  \caption{Looking up the shooting gallery. Looking into the bergschrund at the bottom of the couloir.}
\end{photopair}

\begin{wrapphoto}{r}
  \includegraphics[width=\linewidth]{img/guillaumet/descent_5.jpeg}
  \caption{Looking up at the Brenner route from Paso Guillaumet.}
\end{wrapphoto}
Eventually, we made it out of the bergshrund and back on to moderate snow
slopes. We quickly moved down these and ended up at Paso Guillaumet. We took a
short break here to dry out, and look for our camp mates that were climbing the
Brenner route. We spotted them up on the crux pitch and it looked like they
were starting to bail, a smart decision at that time in the day. We then headed
back towards camp, hiking down skiers right of our approach trail. This hike
was uneventful, except for the relatively large point-release avalanche we saw
that completely wiped out our approach trail! Yikes. Finally, around 2pm, we
made it back to our tent.

Back at camp we re-filled water and sat around for a bit, letting our socks and
boots dry pout in the sun. Ramiro offered to stay at camp for another night to
rest and hike out the next day, but since it was only early afternoon, I decided
we may as well just head down now. A hot meal in El Chalten sounded much better
than another dehydrated meal. So after about an hour, we broke down camp and
re-packed our bags, ready to hike back to town.

\begin{wrapphoto}{l}
  \includegraphics[width=\linewidth]{img/guillaumet/descent_6.jpeg}
  \caption{French fries with cheese, green onion and pancetta.}
\end{wrapphoto}
We made quick work of the descent, all thanks to Ramiro. It was all I could do
to keep pace with him without burning out. Going down the first half was
significantly easier. The long, grinding uphill became a quad burning downhill,
but gravity made my pack feel lighter. Soon we were back at Piedra del Fraille
and we stopped to chat with the owners, and I got a piece of almond cake and a
coke. Although it was a little dry, it was absolutely delicious. After that
digested we continued the walk out. This part was much harder than I remember.
The fatigue was starting to catch up to me and since the trail was completely
flat, it was just as much of an effort as it was coming in. Despite the fatigue,
I was still buzzing with excitement from the climb, and the fact I was in such a
beautiful place. After a few hours we were back close to Rio Electrico and
eventually the car. After the quick drive back to town I found a nearby
restaurant. I ordered some french fires with cheese, green onion and pancetta,
and of course a beer to wash it down. A hearty way to end a hearty trip.

\section{Final Thoughts}

Without any doubt, this was the coolest climbing experience of my life.
Unsurprisingly, Patagonia delivered exactly what I have always heard. The most
striking ridge line in the world, pristine white granite with genuinely good
climbing, and a full value alpine experience (though I didn't experience too
much of the famous wind, thankfully). Despite some of my internal conflicts with
hiring a guide, Ramiro was an awesome partner, a great professional guide, and
someone I genuinely hope to rope up with again. While I wish I could have
tackled an objective like this with a friend or solo, I am very glad I went for
it anyways and swallowed my pride with a guide, it almost certainly wouldn't
have happened without that decision. And now there is a great goal to strive for
on my next trip!

\end{document}
